\begin{enumerate}
\item Tham số hóa tọa độ điểm cần tìm. Nếu điểm thuộc đường thẳng có PTTS thì điểm có tọa độ dạng $(x_0+at; y_0+bt)$.
\item Dùng dữ kiện bài toán (cách đều, khoảng cách, diện tích...) lập phương trình ẩn $t$.
\item Giải phương trình tìm $t$, thay ngược lại để kết luận tọa độ điểm.
\end{enumerate}
Lưu ý: Nếu điểm thuộc đường thẳng dạng PTTQ $ax+by+c=0$, có thể rút $y$ theo $x$ để gọi tọa độ, ví dụ điểm có hoành độ $m$ thì tung độ là $\frac{-am-c}{b}$.
